\documentclass[10pt, aspectratio=169]{beamer}
\usepackage{amsmath, amssymb, amsthm}
\usepackage{xcolor}

\usetheme{Madrid}
\usecolortheme{whale}

% Define custom theorem environments
\newtheorem{intuition}{Intuition}
\newtheorem{motivation}{Motivation}

\title[Phase 3: Optimizations]{Phase 3: Variations and Optimizations}
\subtitle{Fully Homomorphic Encryption over the Integers}
\author{Paper Presentation}
\date{\today}

\begin{document}

\frame{\titlepage}

\begin{frame}{Outline: Phase 3}
    \tableofcontents
\end{frame}

\section{The Practicality Gap}
\begin{frame}{The Practicality Gap}
    \begin{motivation}
    We have established the mathematical engine for our Somewhat Homomorphic Encryption (SHE) scheme. However, before we can secure it or bootstrap it, we face two massive engineering realities:
    \begin{enumerate}
        \item \textbf{Memory Explosion:} Ciphertexts grow exponentially in physical bit-size during evaluation.
        \item \textbf{Bandwidth Choke:} Even fresh ciphertexts are impractically large ($\tilde{\mathcal{O}}(\lambda^5)$ bits).
    \end{enumerate}
    \end{motivation}
    \pause
    
    \vspace{0.3cm}
    To make this scheme practically viable, we must introduce two structural optimizations:
    \begin{itemize}
        \item \textit{In-flight:} Modular reduction during evaluation.
        \item \textit{Post-flight:} Ciphertext compression via group exponentiation.
    \end{itemize}
\end{frame}

\section{Optimization 1: Modular Reduction During Evaluate}
\begin{frame}{The Size Explosion Problem}
    Recall our encryption algorithm: $c = \left[ m + 2r + 2\sum_{i \in S} x_i \right]_{x_0}$.
    \pause
    
    \vspace{0.3cm}
    \textbf{The Problem:}
    During $\text{Encrypt}$, we easily keep the ciphertext size bounded to $\gamma$ bits by reducing modulo $x_0$. 
    
    However, during $\text{Evaluate}$, we \textbf{cannot} simply reduce modulo $x_0$ after a multiplication. 
    \pause
    
    \vspace{0.3cm}
    \textbf{Why?} 
    After one multiplication, the ciphertext squares in size, becoming vastly larger than $x_0$. A standard modulo $x_0$ reduction would subtract a massive multiple of $x_0$. Because $x_0$ contains its own noise ($r_0$), subtracting a massive multiple of $x_0$ injects a catastrophic amount of noise into the ciphertext, instantly destroying the plaintext.
\end{frame}

\begin{frame}{The Solution: A Ladder of Zeroes}
    \textbf{The Fix:} We augment the public key with a "ladder" of increasingly massive encryptions of zero.
    \pause
    
    \vspace{0.2cm}
    For $i = 0, \dots, \gamma$, we generate new elements $x'_i$:
    $$q'_i \leftarrow \mathbb{Z} \cap [2^{\gamma+i-1}/p, 2^{\gamma+i}/p)$$
    $$x'_i = 2(q'_i \cdot p + r'_i)$$
    where $r'_i$ is drawn from the standard noise distribution $(-2^\rho, 2^\rho)$.
    \pause
    
    \vspace{0.3cm}
    \begin{intuition}[The Step-Down Ladder]
    These $x'_i$ values are astronomically large, scaling up to $2^{2\gamma}$. 
    
    During evaluation, when a ciphertext exceeds $2^\gamma$, we do not reduce by $x_0$. Instead, we reduce modulo the largest $x'_\gamma$, then modulo $x'_{\gamma-1}$, and so on, stepping down the ladder until the ciphertext is back to $\gamma$ bits.
    \end{intuition}
    \pause
    
    \textbf{Result:} Because we only ever subtract small multiples of these $x'_i$ elements at each step, the accumulated error remains strictly bounded!
\end{frame}

\section{Optimization 2: Ciphertext Compression}
\begin{frame}{The Bandwidth Problem}
    Even with the ladder optimization keeping ciphertexts at $\gamma$ bits, $\gamma$ is still massively large—specifically $\tilde{\mathcal{O}}(\lambda^5)$ bits. Transmitting this over a network is brutally inefficient.
    \pause
    
    \vspace{0.3cm}
    \textbf{The Goal:} Compress the final output ciphertext to the size of a standard RSA modulus \textit{after} the cloud server has finished all homomorphic evaluations.
\end{frame}

\begin{frame}{Compression via Group Exponentiation}
    \textbf{The Mechanism:}
    We supplement the public key with the description of a cyclic group $G = \langle g \rangle$ whose order is a multiple of the secret key $p$.
    \pause
    
    \vspace{0.3cm}
    \textbf{The Compression Step:}
    Instead of sending the massive integer $c$ back to the user, the server computes and transmits:
    $$c^* = g^c \pmod N$$
    \pause
    
    \vspace{0.3cm}
    \textbf{The Decryption Step:}
    Because $c^*$ is an element of the group, the user decrypts by computing the Discrete Logarithm (DL):
    $$m = \left( DL_g(c^*) \pmod p \right) \pmod 2$$
\end{frame}

\begin{frame}{Trade-offs of Compression}
    How is computing the Discrete Logarithm efficient for the user?
    \pause
    
    \vspace{0.3cm}
    \textbf{The Secret Key Modification:}
    We must generate the secret key $p$ as a \textbf{smooth} number (a product of many small distinct primes). This allows the user to efficiently compute the discrete log using the Pohlig-Hellman algorithm or baby-step giant-step method.
    \pause
    
    \vspace{0.3cm}
    \begin{block}{The Ultimate Trade-off}
    \begin{itemize}
        \item \textbf{Pros:} The ciphertext is compressed from $\tilde{\mathcal{O}}(\lambda^5)$ bits down to the size of an RSA modulus ($\sim 2048$ bits).
        \item \textbf{Cons:} The compressed ciphertext $c^*$ resides in the group $G$. Because we do not have an FHE scheme over this group, \textbf{we completely lose all homomorphic properties.} 
    \end{itemize}
    \end{block}
    \pause
    
    \vspace{0.2cm}
    \textit{Conclusion:} This is strictly a post-flight optimization, applied only when computation is complete and data is ready for transit.
\end{frame}

\section{References}
\begin{frame}{References}
    \begin{thebibliography}{9}
        \bibitem{DGHV10}
        Marten van Dijk, Craig Gentry, Shai Halevi, and Vinod Vaikuntanathan.
        \textit{Fully Homomorphic Encryption over the Integers}. 
        Eurocrypt, 2010.
        
        \bibitem{Cachin99}
        C. Cachin, S. Micali, and M. Stadler.
        \textit{Computationally private information retrieval with polylogarithmic communication}.
        EUROCRYPT, 1999. (For the $\phi$-hiding assumption).
    \end{thebibliography}
\end{frame}

\end{document}